% =====================================================================
% Perceiver & Perceiver IO -- Presentazione Beamer
% Obiettivo: PDF a slide "tipo PowerPoint" generato da LaTeX
% Ordine: Note presentazione.pptx
% Fonti: presentation_v2_eng.pdf/tex, presentation_v2_ita.tex, paper Perceiver
% =====================================================================
\documentclass[aspectratio=169,10pt]{beamer}

\usepackage[T1]{fontenc}
\usepackage[utf8]{inputenc}
\usepackage[italian]{babel}
\usepackage{lmodern}
\usepackage{amsmath,amssymb,booktabs,array}
\usepackage{tikz}
\usetikzlibrary{arrows.meta,positioning,calc,shapes.geometric}

\usetheme{Madrid}
\usecolortheme{default}

\definecolor{DeepBlue}{HTML}{123C69}
\definecolor{Ink}{HTML}{1D2630}
\definecolor{SoftBlue}{HTML}{E8F1FA}
\definecolor{SoftGreen}{HTML}{E9F7EF}
\definecolor{SoftAmber}{HTML}{FFF4D8}
\definecolor{AccentOrange}{HTML}{D97706}
\definecolor{AccentGreen}{HTML}{168056}
\definecolor{AccentRed}{HTML}{B42318}

\setbeamercolor{structure}{fg=DeepBlue}
\setbeamercolor{title}{fg=white,bg=DeepBlue}
\setbeamercolor{frametitle}{fg=white,bg=DeepBlue}
\setbeamercolor{normal text}{fg=Ink,bg=white}
\setbeamercolor{block title}{bg=DeepBlue,fg=white}
\setbeamercolor{block body}{bg=SoftBlue,fg=Ink}
\setbeamercolor{block title example}{bg=AccentGreen,fg=white}
\setbeamercolor{block body example}{bg=SoftGreen,fg=Ink}
\setbeamercolor{block title alerted}{bg=AccentRed,fg=white}
\setbeamercolor{block body alerted}{bg=AccentRed!8,fg=Ink}
\setbeamercolor{alerted text}{fg=AccentOrange}
\setbeamertemplate{navigation symbols}{}
\setbeamertemplate{itemize item}{\color{DeepBlue}$\blacktriangleright$}
\setbeamertemplate{itemize subitem}{\color{AccentOrange}$\bullet$}
\setbeamertemplate{footline}{
  \leavevmode%
  \hbox{%
    \begin{beamercolorbox}[wd=.42\paperwidth,ht=2.4ex,dp=1ex,center]{author in head/foot}%
      \usebeamerfont{author in head/foot}\scriptsize F.Gigli C.Vignoli
    \end{beamercolorbox}%
    \begin{beamercolorbox}[wd=.40\paperwidth,ht=2.4ex,dp=1ex,center]{title in head/foot}%
      \usebeamerfont{title in head/foot}\scriptsize Perceiver \& Perceiver IO
    \end{beamercolorbox}%
    \begin{beamercolorbox}[wd=.18\paperwidth,ht=2.4ex,dp=1ex,right]{date in head/foot}%
      \usebeamerfont{date in head/foot}\scriptsize\insertframenumber{}/\inserttotalframenumber\hspace*{1.5ex}
    \end{beamercolorbox}%
  }%
}

\newcommand{\R}{\mathbb{R}}
\newcommand{\bigO}{\mathcal{O}}
\newcommand{\softmax}{\operatorname{softmax}}
\newcommand{\good}[1]{\textcolor{AccentGreen}{\textbf{#1}}}
\newcommand{\warn}[1]{\textcolor{AccentOrange}{\textbf{#1}}}
\newcommand{\bad}[1]{\textcolor{AccentRed}{\textbf{#1}}}

\title[Perceiver \& Perceiver IO]{Perceiver \& Perceiver IO}
\subtitle{Implementazione e analisi sperimentale}
\author[F.Gigli C.Vignoli]{Francesco Gigli e Corso Vignoli}
\institute[UNIFI]{Università degli Studi di Firenze\\Corso B031278 -- Deep Learning}
\date{}

\begin{document}

% ---------------------------------------------------------------------
\begin{frame}
  \titlepage
\end{frame}

% ---------------------------------------------------------------------
\begin{frame}{Agenda}
  \begin{enumerate}
    \item Introduzione e motivazione
    \item Architettura
    \item Esperimenti e risultati
    \item Criticità e discussione
  \end{enumerate}
\end{frame}

\section{Introduzione e motivazione}

% ---------------------------------------------------------------------
\begin{frame}{Il problema: architetture specifiche per modalità}
\small
\begin{columns}[T]
\begin{column}{0.52\textwidth}
\textbf{Approccio tradizionale}
\begin{itemize}
  \item Immagini $\rightarrow$ CNN: ResNet, EfficientNet
  \item Testo $\rightarrow$ Transformer: BERT, GPT
  \item Audio $\rightarrow$ CNN 1D / RNN
  \item Point cloud $\rightarrow$ PointNet, DGCNN
\end{itemize}

\begin{alertblock}{Problema}
Ogni modalità richiede assunzioni e architetture diverse: località, griglia, ordine, geometria.
\end{alertblock}
\end{column}
\begin{column}{0.45\textwidth}
\centering
\begin{tikzpicture}[node distance=0.45cm,
  box/.style={draw,rounded corners=2pt,minimum width=2.05cm,minimum height=0.48cm,font=\scriptsize\bfseries},
  arr/.style={-Stealth,thick,AccentRed}]
\node[box,fill=SoftBlue] (img) {Immagini};
\node[box,fill=SoftBlue,below=of img] (txt) {Testo};
\node[box,fill=SoftBlue,below=of txt] (aud) {Audio};
\node[box,fill=SoftBlue,below=of aud] (pc) {Point cloud};
\node[box,fill=AccentRed!13,right=1.6cm of img] (cnn) {CNN};
\node[box,fill=AccentRed!13,right=1.6cm of txt] (tra) {Transformer};
\node[box,fill=AccentRed!13,right=1.6cm of aud] (rnn) {CNN/RNN};
\node[box,fill=AccentRed!13,right=1.6cm of pc] (pnt) {PointNet};
\draw[arr] (img) -- (cnn);
\draw[arr] (txt) -- (tra);
\draw[arr] (aud) -- (rnn);
\draw[arr] (pc) -- (pnt);
\end{tikzpicture}

\vspace{0.4em}
\textbf{Domanda:}\\
si può usare una sola architettura per dati molto diversi?
\end{column}
\end{columns}
\end{frame}

% ---------------------------------------------------------------------
\begin{frame}{La soluzione Perceiver}
\small
\begin{columns}[T]
\begin{column}{0.55\textwidth}
\textbf{Idea chiave} (Jaegle et al., ICML 2021):
\begin{enumerate}
  \item introdurre un \alert{latent array} compatto $M \ll N$;
  \item usare cross-attention: latenti = query, input = key/value;
  \item raffinare i latenti con self-attention;
  \item ripetere il blocco con weight sharing.
\end{enumerate}
\end{column}
\begin{column}{0.42\textwidth}
\begin{block}{Complessità}
\[
\text{Transformer: } \bigO(N^2 d)
\]
\[
\text{Perceiver: } \bigO(NMd + M^2Ld)
\]
\textbf{Lineare in $N$} se $M$ è fissato e $M \ll N$.
\end{block}

\begin{exampleblock}{Contributo del progetto}
Implementazione PyTorch from-scratch su immagini, point cloud e testo.
\end{exampleblock}
\end{column}
\end{columns}
\end{frame}

% ---------------------------------------------------------------------
\begin{frame}{Perceiver IO: output strutturati}
\small
\begin{center}
\begin{tikzpicture}[
  block/.style={draw,thick,rounded corners=2pt,minimum width=2.4cm,minimum height=0.8cm,align=center,font=\small\bfseries},
  arr/.style={-Stealth,very thick}]
\node[block,fill=SoftBlue] (read) {Read\\cross-attn};
\node[block,fill=SoftAmber,right=1.2cm of read] (process) {Process\\latent self-attn};
\node[block,fill=SoftGreen,right=1.2cm of process] (write) {Write\\output queries};
\draw[arr,DeepBlue] (read) -- (process);
\draw[arr,AccentGreen] (process) -- (write);
\node[below=0.5cm of read,font=\scriptsize] {input $\rightarrow$ latenti};
\node[below=0.5cm of process,font=\scriptsize] {raffinamento};
\node[below=0.5cm of write,font=\scriptsize] {latenti $\rightarrow$ output};
\end{tikzpicture}
\end{center}

\begin{columns}[T]
\begin{column}{0.48\textwidth}
\textbf{Perceiver base}
\begin{itemize}
  \item encoder general-purpose;
  \item output semplice, adatto a classificazione;
  \item readout tramite pooling/head.
\end{itemize}
\end{column}
\begin{column}{0.48\textwidth}
\textbf{Perceiver IO}
\begin{itemize}
  \item output queries task-specific;
  \item output di dimensione arbitraria;
  \item classificazione, MLM, optical flow, multimodale.
\end{itemize}
\end{column}
\end{columns}
\end{frame}

\section{Architettura}

% ---------------------------------------------------------------------
\begin{frame}[plain]
\centering
\vfill
{\Huge\color{DeepBlue}\textbf{Architettura}}
\vspace{1em}

{\large Dal dato grezzo allo spazio latente}
\vfill
\end{frame}

% ---------------------------------------------------------------------
\begin{frame}{Panoramica architetturale}
\small
\begin{center}
\begin{tikzpicture}[
  node distance=0.75cm,
  block/.style={draw,rounded corners=2pt,thick,align=center,minimum width=2.05cm,minimum height=0.75cm,font=\scriptsize\bfseries},
  arr/.style={-Stealth,thick}]
\node[block,fill=SoftBlue] (input) {Input array\\$N \times C$};
\node[block,fill=SoftAmber,right=of input] (pe) {Fourier PE\\$N \times C_{\text{tot}}$};
\node[block,fill=SoftGreen,right=of pe] (cross) {Cross-attn\\$M \times N$};
\node[block,fill=SoftBlue,right=of cross] (latent) {Latent\\Transformer};
\node[block,fill=SoftGreen,right=of latent] (out) {Output};
\draw[arr] (input) -- (pe);
\draw[arr] (pe) -- (cross);
\draw[arr] (cross) -- (latent);
\draw[arr] (latent) -- (out);
\end{tikzpicture}
\end{center}

\begin{itemize}
  \item \textbf{Input array}: sequenza generica di pixel, punti 3D o byte.
  \item \textbf{Latent array}: bottleneck appreso $M \times D$, con $M \ll N$.
  \item \textbf{Weight sharing}: stesso blocco riapplicato in più iterazioni.
  \item \textbf{Output}: pooling per Perceiver, queries per Perceiver IO.
\end{itemize}
\end{frame}

% ---------------------------------------------------------------------
\begin{frame}{Cross-attention: il meccanismo chiave}
\small
\begin{block}{Formulazione}
\[
Q = X_{\text{latent}} W_Q \in \R^{M \times d_k}
\]
\[
K = X_{\text{input}} W_K \in \R^{N \times d_k}, \qquad
V = X_{\text{input}} W_V \in \R^{N \times d_v}
\]
\[
\operatorname{CrossAttn}(Q,K,V)=
\softmax\!\left(\frac{QK^\top}{\sqrt{d_k}}\right)V
\]
\end{block}

\begin{itemize}
  \item I latenti fanno le \textbf{query}: decidono cosa leggere.
  \item L'input fornisce \textbf{key/value}: descrizione e contenuto.
  \item La matrice di attenzione è $M \times N$, non $N \times N$.
\end{itemize}
\end{frame}

% ---------------------------------------------------------------------
\begin{frame}{Latent Transformer e weight sharing}
\small
\begin{columns}[T]
\begin{column}{0.52\textwidth}
\textbf{Self-attention tra latenti}
\begin{itemize}
  \item $Q=K=V=X_{\text{latent}}$
  \item costo $\bigO(M^2)$
  \item pre-norm LayerNorm + residual
  \item feed-forward con espansione $4\times$
  \item GELU/GEGLU come non linearità
\end{itemize}
\end{column}
\begin{column}{0.45\textwidth}
\begin{block}{Weight sharing}
Un singolo blocco viene applicato più volte:
\[
Z^{t+1}=F_\theta(Z^t, X)
\]
È simile a un RNN nello spazio latente.
\end{block}
\begin{exampleblock}{Nel progetto}
Shared: $\sim3.35$M parametri\\
Unshared: $\sim11.0$M parametri\\
Riduzione: circa $3\times$
\end{exampleblock}
\end{column}
\end{columns}
\end{frame}

% ---------------------------------------------------------------------
\begin{frame}{Fourier Positional Encoding}
\small
\begin{columns}[T]
\begin{column}{0.50\textwidth}
\textbf{Perché serve?}
\begin{itemize}
  \item L'attention è permutation-equivariant.
  \item Senza posizione, il modello non sa \emph{dove} sono pixel/punti/token.
  \item L'immagine diventa quasi un insieme non ordinato di patch.
\end{itemize}
\end{column}
\begin{column}{0.47\textwidth}
\begin{block}{Fourier features}
\[
\gamma(x)=
\left[
\sin(2\pi f_k x_i),
\cos(2\pi f_k x_i)
\right]_{k,i}
\]
Frequenze multi-scala concatenate ai canali originali.
\end{block}
\end{column}
\end{columns}

\begin{alertblock}{Risultato critico su CIFAR-10}
Con Fourier PE: \textbf{72.23\%}. Senza PE: \textbf{35.41\%}. Drop: \textbf{-36.82 pp}.
\end{alertblock}
\end{frame}

% ---------------------------------------------------------------------
\begin{frame}{Perceiver IO: decoder con output queries}
\small
\begin{block}{Decoder cross-attention}
\[
Q_{\text{dec}} = Q_{\text{out}} W_Q,\qquad
K_{\text{dec}} = Z W_K,\qquad
V_{\text{dec}} = Z W_V
\]
\[
Y =
\softmax\!\left(\frac{Q_{\text{dec}}K_{\text{dec}}^\top}{\sqrt{d_k}}\right)V_{\text{dec}}
\in \R^{O \times E}
\]
\end{block}

\begin{itemize}
  \item $O$ = numero di output queries, dipendente dal task.
  \item $O$ è indipendente da lunghezza input $N$ e numero latenti $M$.
  \item Una sola cross-attention di decoder può bastare per leggere i latenti.
\end{itemize}
\end{frame}

% ---------------------------------------------------------------------
\begin{frame}{Design delle output queries per task}
\small
\begin{center}
\begin{tabular}{lll}
\toprule
\textbf{Task} & \textbf{Query} & \textbf{Output} \\
\midrule
Classificazione & $O=1$ & logits di classe \\
Optical flow & $O=H \times W$ & flow per pixel \\
Masked LM & posizioni mascherate & logits vocab \\
Multimodale & query/task embedding & output misti \\
\bottomrule
\end{tabular}
\end{center}

\vspace{0.5em}
\begin{exampleblock}{Nel progetto}
\begin{itemize}
  \item CIFAR-10 e ModelNet40: una query / aggregazione per classificazione.
  \item WikiText-103: una query per posizione mascherata, vocab byte-level $=256$.
\end{itemize}
\end{exampleblock}
\end{frame}

% ---------------------------------------------------------------------
\begin{frame}{Analisi della complessità computazionale}
\small
\begin{columns}[T]
\begin{column}{0.48\textwidth}
\begin{block}{Modelli}
\[
\text{Transformer: } \bigO(N^2 d L)
\]
\[
\text{Perceiver: } \bigO((NM + M^2L)d)
\]
\[
\text{Perceiver IO: } \bigO(((N+O)M + M^2L)d)
\]
\end{block}
\end{column}
\begin{column}{0.48\textwidth}
\begin{exampleblock}{Esempio CIFAR-10}
\[
N=1024,\quad M=96,\quad L=4
\]
\[
N^2 = 1\,048\,576
\]
\[
NM+M^2L = 135\,168
\]
\textbf{$\sim7.8\times$ meno interazioni}
\end{exampleblock}
\end{column}
\end{columns}

\begin{alertblock}{Scaling}
Su immagini $224 \times 224$, il vantaggio diventa nell'ordine di $\sim100\times$.
\end{alertblock}
\end{frame}

% ---------------------------------------------------------------------
\begin{frame}{GELU e ottimizzatore LAMB}
\small
\begin{columns}[T]
\begin{column}{0.48\textwidth}
\begin{block}{GELU}
\[
\operatorname{GELU}(x)=x\Phi(x)\approx x\sigma(1.702x)
\]
\begin{itemize}
  \item liscia;
  \item standard nei Transformer;
  \item stabile in MLP profonde.
\end{itemize}
\end{block}
\end{column}
\begin{column}{0.48\textwidth}
\begin{block}{LAMB}
\[
r=\frac{\lVert w\rVert}{\lVert w+\eta \hat u\rVert}
\]
\begin{itemize}
  \item trust ratio per layer;
  \item utile per batch grandi;
  \item usato con warmup + cosine annealing.
\end{itemize}
\end{block}
\end{column}
\end{columns}
\end{frame}

% ---------------------------------------------------------------------
\begin{frame}{Implementazione: struttura del codice}
\small
\begin{columns}[T]
\begin{column}{0.52\textwidth}
\textbf{Componenti principali}
\begin{itemize}
  \item \texttt{model.py}, \texttt{perceiver\_io.py}
  \item \texttt{attention.py}
  \item \texttt{positional\_encoding.py}
  \item \texttt{lamb\_optimizer.py}
\end{itemize}
\end{column}
\begin{column}{0.45\textwidth}
\textbf{Esperimenti}
\begin{itemize}
  \item CIFAR-10
  \item ModelNet40
  \item WikiText-103
  \item GLUE
\end{itemize}
\end{column}
\end{columns}

\begin{exampleblock}{Numeri}
Core model $\sim600$ righe; data modules $\sim500$; esperimenti $\sim800$; totale $\sim2500$ righe.
\end{exampleblock}
\end{frame}

\section{Esperimenti e risultati}

% ---------------------------------------------------------------------
\begin{frame}[plain]
\centering
\vfill
{\Huge\color{DeepBlue}\textbf{Esperimenti e risultati}}
\vspace{1em}

{\large Immagini, point cloud e testo}
\vfill
\end{frame}

% ---------------------------------------------------------------------
\begin{frame}{CIFAR-10: setup sperimentale}
\small
\begin{columns}[T]
\begin{column}{0.50\textwidth}
\begin{tabular}{ll}
\toprule
\textbf{Parametro} & \textbf{Valore} \\
\midrule
Latenti & $96 \times 384$ \\
Cross-attn stages & $4$ \\
Transformer blocks & $4$ shared \\
Heads & $3$ \\
Patch size & $2 \times 2$ \\
Dropout & $0.2$ \\
\bottomrule
\end{tabular}
\end{column}
\begin{column}{0.47\textwidth}
\begin{tabular}{ll}
\toprule
\textbf{Training} & \textbf{Valore} \\
\midrule
Optimizer & LAMB \\
LR & $0.004$ \\
Batch size & $64$ \\
Epochs & $120$ \\
Parametri & $\sim3.35$M \\
\bottomrule
\end{tabular}
\end{column}
\end{columns}

\begin{block}{Piano ablation}
Positional encoding, weight sharing, permutazione pixel, Perceiver vs Perceiver IO.
\end{block}
\end{frame}

% ---------------------------------------------------------------------
\begin{frame}{CIFAR-10: risultati ablation}
\small
\begin{center}
\begin{tabular}{lc}
\toprule
\textbf{Configurazione} & \textbf{Accuracy} \\
\midrule
Perceiver IO & \good{78.20\%} \\
No weight sharing & 73.61\% \\
Fourier control & 72.23\% \\
Baseline Fourier & 69.69\% \\
WS control & 68.70\% \\
Fourier permuted & 62.31\% \\
Learned PE permuted & 55.12\% \\
RGB only, no PE & \bad{35.41\%} \\
\bottomrule
\end{tabular}
\end{center}

\begin{exampleblock}{Osservazioni}
PE è essenziale; Perceiver IO migliora il readout; weight sharing resta competitivo con meno parametri.
\end{exampleblock}
\end{frame}

% ---------------------------------------------------------------------
\begin{frame}{CIFAR-10: impatto del positional encoding}
\small
\begin{columns}[T]
\begin{column}{0.48\textwidth}
\begin{block}{Risultato}
\[
\text{Fourier PE}=72.23\%
\]
\[
\text{RGB only}=35.41\%
\]
\[
\Delta=-36.82\text{ pp}
\]
\end{block}
\end{column}
\begin{column}{0.48\textwidth}
\textbf{Perché succede?}
\begin{itemize}
  \item senza PE l'attention non vede la posizione;
  \item l'immagine diventa una bag of patches;
  \item la struttura 2D viene persa;
  \item conferma l'ipotesi del paper.
\end{itemize}
\end{column}
\end{columns}

\begin{alertblock}{Takeaway}
Il positional encoding è il componente più importante della riproduzione.
\end{alertblock}
\end{frame}

% ---------------------------------------------------------------------
\begin{frame}{CIFAR-10: weight sharing}
\small
\begin{center}
\begin{tabular}{lcc}
\toprule
\textbf{Modello} & \textbf{Parametri} & \textbf{Accuracy} \\
\midrule
Shared baseline & 3.35M & 69.69\% \\
WS control & 3.35M & 68.70\% \\
No sharing & 11.0M & 73.61\% \\
\bottomrule
\end{tabular}
\end{center}

\begin{columns}[T]
\begin{column}{0.48\textwidth}
\begin{exampleblock}{Vantaggio}
Circa $3\times$ meno parametri, memoria minore, regolarizzazione implicita.
\end{exampleblock}
\end{column}
\begin{column}{0.48\textwidth}
\begin{alertblock}{Costo}
No sharing guadagna $+3.92$ pp ma usa molti più parametri.
\end{alertblock}
\end{column}
\end{columns}
\end{frame}

% ---------------------------------------------------------------------
\begin{frame}{CIFAR-10: robustezza alla permutazione}
\small
\begin{center}
\begin{tabular}{lcc}
\toprule
\textbf{Configurazione} & \textbf{Accuracy} & \textbf{Drop} \\
\midrule
Fourier PE normale & 72.23\% & -- \\
Fourier PE permutato & 62.31\% & -9.92 pp \\
Learned PE permutato & 55.12\% & -17.11 pp \\
\bottomrule
\end{tabular}
\end{center}

\begin{block}{Analisi}
Fourier PE conserva meglio informazione multi-scala; learned PE tende a memorizzare posizioni assolute.
\end{block}
\end{frame}

% ---------------------------------------------------------------------
\begin{frame}{CIFAR-10: Perceiver vs Perceiver IO}
\small
\begin{center}
\begin{tabular}{lcc}
\toprule
\textbf{Modello} & \textbf{Readout} & \textbf{Accuracy} \\
\midrule
Perceiver & mean pooling / head & 72.23\% \\
Perceiver IO & output query decoder & \good{78.20\%} \\
\bottomrule
\end{tabular}
\end{center}

\begin{exampleblock}{Perché Perceiver IO migliora?}
\begin{itemize}
  \item la query di output è un aggregatore appreso;
  \item il decoder può selezionare latenti rilevanti;
  \item aggiunge capacità senza cambiare l'encoder.
\end{itemize}
\end{exampleblock}
\end{frame}

% ---------------------------------------------------------------------
\begin{frame}{ModelNet40: setup e risultati}
\small
\begin{columns}[T]
\begin{column}{0.50\textwidth}
\begin{tabular}{ll}
\toprule
\textbf{Parametro} & \textbf{Valore} \\
\midrule
Latenti & $128 \times 512$ \\
Cross-attn stages & $2$ \\
Transformer blocks & $6$ shared \\
Punti/oggetto & $2048$ \\
Optimizer & LAMB \\
Epochs & $200$ \\
Batch size & $128$ \\
\bottomrule
\end{tabular}
\end{column}
\begin{column}{0.47\textwidth}
\begin{tabular}{lc}
\toprule
\textbf{Setup} & \textbf{Accuracy} \\
\midrule
Baseline scale & \good{84.16\%} \\
+ rotation & 83.06\% \\
+ translation & 82.90\% \\
Paper originale & 85.7\% \\
\bottomrule
\end{tabular}
\end{column}
\end{columns}

\begin{block}{Risultato}
Gap di soli \textbf{1.54 pp} dal paper originale: riproduzione molto vicina.
\end{block}
\end{frame}

% ---------------------------------------------------------------------
\begin{frame}{ModelNet40: effetto della data augmentation}
\small
\begin{center}
\begin{tabular}{lcc}
\toprule
\textbf{Augmentation} & \textbf{Accuracy} & \textbf{Delta} \\
\midrule
Scale only & 84.16\% & -- \\
Scale + rotation & 83.06\% & -1.10 pp \\
Scale + translation & 82.90\% & -1.26 pp \\
\bottomrule
\end{tabular}
\end{center}

\begin{alertblock}{Perché può peggiorare?}
\begin{itemize}
  \item capacità limitata: 128 latenti;
  \item augmentation aggressiva richiede più training;
  \item Fourier PE su coordinate assolute è sensibile a rotazioni/traslazioni.
\end{itemize}
\end{alertblock}
\end{frame}

% ---------------------------------------------------------------------
\begin{frame}{WikiText-103: pre-training MLM byte-level}
\small
\begin{columns}[T]
\begin{column}{0.50\textwidth}
\begin{tabular}{ll}
\toprule
\textbf{Parametro} & \textbf{Valore} \\
\midrule
Tokenizzazione & byte-level UTF-8 \\
Vocabolario & 256 \\
Sequence length & 1024 \\
Mask probability & 15\% \\
Modello & Perceiver IO \\
Parametri & $\sim10$M \\
Epochs & 50 \\
\bottomrule
\end{tabular}
\end{column}
\begin{column}{0.47\textwidth}
\begin{exampleblock}{Perché byte-level?}
\begin{itemize}
  \item nessun BPE/WordPiece;
  \item input davvero language-agnostic;
  \item vocabolario fisso e piccolo.
\end{itemize}
\end{exampleblock}
\begin{alertblock}{Trade-off}
Sequenze più lunghe per rappresentare lo stesso testo.
\end{alertblock}
\end{column}
\end{columns}
\end{frame}

% ---------------------------------------------------------------------
\begin{frame}{GLUE: fine-tuning su 8 task}
\small
\begin{columns}[T]
\begin{column}{0.48\textwidth}
\textbf{Pipeline}
\begin{enumerate}
  \item Pre-training MLM su WikiText-103
  \item Trasferimento encoder
  \item Nuova query di classificazione
  \item Fine-tuning GLUE
\end{enumerate}
\end{column}
\begin{column}{0.49\textwidth}
\textbf{Task GLUE}
\begin{itemize}
  \item Singola frase: CoLA, SST-2
  \item Similarità: MRPC, STS-B, QQP
  \item Inferenza: MNLI, QNLI, RTE
\end{itemize}
\end{column}
\end{columns}

\begin{block}{Setup}
30 epoche, learning rate $5 \times 10^{-4}$, input byte-level.
\end{block}
\end{frame}

% ---------------------------------------------------------------------
\begin{frame}{GLUE: analisi per task}
\small
\begin{columns}[T]
\begin{column}{0.50\textwidth}
\begin{exampleblock}{Osservazioni}
\begin{itemize}
  \item sopra majority baseline su molti task;
  \item SST-2 relativamente forte;
  \item transfer MLM $\rightarrow$ GLUE utile.
\end{itemize}
\end{exampleblock}
\end{column}
\begin{column}{0.47\textwidth}
\begin{alertblock}{Limitazioni}
\begin{itemize}
  \item byte-level = sequenze lunghe;
  \item modello $\sim10$M vs BERT base 110M;
  \item niente semantica subword.
\end{itemize}
\end{alertblock}
\end{column}
\end{columns}
\end{frame}

% ---------------------------------------------------------------------
\begin{frame}{Convergenza: tutti gli esperimenti}
\small
\begin{center}
\begin{tabular}{lc}
\toprule
\textbf{Modalità} & \textbf{Convergenza indicativa} \\
\midrule
Point cloud & $\sim50$ epoche \\
GLUE fine-tuning & 15--30 epoche \\
CIFAR-10 & $\sim120$ epoche \\
\bottomrule
\end{tabular}
\end{center}

\begin{itemize}
  \item Point cloud: input a bassa dimensione, convergenza più rapida.
  \item GLUE: il pre-training accelera l'adattamento.
  \item CIFAR-10: pixel/patch grezzi richiedono più iterazioni.
\end{itemize}
\end{frame}

% ---------------------------------------------------------------------
\begin{frame}{Attention maps: analisi visiva}
\small
\begin{columns}[T]
\begin{column}{0.48\textwidth}
\begin{exampleblock}{Con Fourier PE}
\begin{itemize}
  \item pattern spazialmente coerenti;
  \item latenti specializzati;
  \item attenzione meno entropica.
\end{itemize}
\end{exampleblock}
\end{column}
\begin{column}{0.48\textwidth}
\begin{alertblock}{Senza PE}
\begin{itemize}
  \item pattern meno strutturati;
  \item entropia più alta;
  \item difficile separare regioni informative.
\end{itemize}
\end{alertblock}
\end{column}
\end{columns}

\begin{block}{Interpretazione}
Le attention maps confermano qualitativamente il risultato quantitativo: la posizione guida la lettura dell'input.
\end{block}
\end{frame}

% ---------------------------------------------------------------------
\begin{frame}{Risultati complessivi}
\small
\begin{center}
\begin{tabular}{lll}
\toprule
\textbf{Modalità} & \textbf{Risultato} & \textbf{Nota} \\
\midrule
Immagini & Perceiver IO 78.20\% & migliore CIFAR-10 \\
Point cloud & ModelNet40 84.16\% & gap 1.54 pp dal paper \\
Testo & MLM + 8 task GLUE & input byte-level \\
\bottomrule
\end{tabular}
\end{center}

\begin{exampleblock}{Sintesi}
\begin{itemize}
  \item $\sim2500$ righe PyTorch from-scratch;
  \item 20+ training run;
  \item 3 modalità: immagini, point cloud, testo;
  \item pipeline riproducibile.
\end{itemize}
\end{exampleblock}
\end{frame}

\section{Criticità e discussione}

% ---------------------------------------------------------------------
\begin{frame}[plain]
\centering
\vfill
{\Huge\color{DeepBlue}\textbf{Criticità e discussione}}
\vspace{1em}

{\large Dove funziona, dove costa, cosa migliorare}
\vfill
\end{frame}

% ---------------------------------------------------------------------
\begin{frame}{Gap rispetto al paper originale}
\small
\begin{center}
\begin{tabular}{lccc}
\toprule
\textbf{Dataset} & \textbf{Paper} & \textbf{Progetto} & \textbf{Gap} \\
\midrule
ModelNet40 & 85.7\% & 84.16\% & \good{1.54 pp} \\
CIFAR-10 / immagini & $>90\%$ & 78.20\% & \bad{$\sim12$ pp} \\
\bottomrule
\end{tabular}
\end{center}

\begin{block}{Interpretazione}
La riproduzione sui point cloud è molto vicina. Il gap sulle immagini dipende soprattutto dal compute disponibile.
\end{block}
\end{frame}

% ---------------------------------------------------------------------
\begin{frame}{Limitazioni hardware e scelte implementative}
\small
\begin{center}
\begin{tabular}{lll}
\toprule
\textbf{Aspetto} & \textbf{Paper Google} & \textbf{Progetto} \\
\midrule
Hardware & 512 TPU v3 cores & 1 GPU consumer \\
Batch size & fino a 1024 & 64--128 \\
Latenti & 256--512 & 96--128 \\
Parametri & 30--100M+ & 3.35--10M \\
\bottomrule
\end{tabular}
\end{center}

\begin{alertblock}{Conseguenze}
LAMB meno efficace con batch piccoli; meno latenti limitano il bottleneck; modello più piccolo = meno capacità.
\end{alertblock}
\end{frame}

% ---------------------------------------------------------------------
\begin{frame}{Possibili miglioramenti}
\small
\begin{columns}[T]
\begin{column}{0.48\textwidth}
\textbf{Architettura}
\begin{itemize}
  \item aumentare $M$ con gradient checkpointing;
  \item cross-attention multi-scala;
  \item positional encoding relativi: RoPE, ALiBi.
\end{itemize}

\textbf{Training}
\begin{itemize}
  \item batch effettivo più grande;
  \item mixed precision FP16/BF16;
  \item CutMix e MixUp.
\end{itemize}
\end{column}
\begin{column}{0.48\textwidth}
\textbf{Dati e analisi}
\begin{itemize}
  \item tokenizzazione subword per NLP;
  \item training multimodale congiunto;
  \item confronto con Linformer, Performer e altri efficient Transformer.
\end{itemize}
\end{column}
\end{columns}
\end{frame}

% ---------------------------------------------------------------------
\begin{frame}{Lezioni apprese}
\small
\begin{columns}[T]
\begin{column}{0.48\textwidth}
\begin{exampleblock}{Cosa funziona}
\begin{itemize}
  \item architettura davvero modality-agnostic;
  \item weight sharing efficiente;
  \item Perceiver IO migliora il readout;
  \item transfer MLM $\rightarrow$ GLUE utile.
\end{itemize}
\end{exampleblock}
\end{column}
\begin{column}{0.48\textwidth}
\begin{alertblock}{Cosa limita}
\begin{itemize}
  \item gap immagini con poco compute;
  \item forte dipendenza dal PE;
  \item LAMB sensibile al batch size;
  \item byte-level = sequenze lunghe.
\end{itemize}
\end{alertblock}
\end{column}
\end{columns}

\begin{block}{Main takeaway}
Il positional encoding è il componente più critico: senza posizione, CIFAR-10 crolla da 72\% a 35\%.
\end{block}
\end{frame}

% ---------------------------------------------------------------------
\begin{frame}{Conclusioni}
\small
\begin{columns}[T]
\begin{column}{0.50\textwidth}
\textbf{Punti di forza}
\begin{itemize}
  \item una sola architettura per immagini, point cloud e testo;
  \item costo lineare nella dimensione input;
  \item meno memoria con weight sharing;
  \item nessuna CNN/RNN specifica per dominio.
\end{itemize}
\end{column}
\begin{column}{0.47\textwidth}
\textbf{Debolezze}
\begin{itemize}
  \item richiede molto compute per battere modelli specializzati;
  \item dipende dal positional encoding;
  \item bottleneck latente può perdere dettagli fini.
\end{itemize}
\end{column}
\end{columns}

\begin{exampleblock}{Messaggio finale}
Il Perceiver è un passo convincente verso modelli di percezione general-purpose.
\end{exampleblock}
\end{frame}

\section{Chiusura}

% ---------------------------------------------------------------------
\begin{frame}{Riferimenti}
\small
\begin{thebibliography}{9}
\bibitem{perceiver} Jaegle et al., \emph{Perceiver: General Perception with Iterative Attention}, ICML 2021.
\bibitem{perceiverio} Jaegle et al., \emph{Perceiver IO: A General Architecture for Structured Inputs \& Outputs}, ICLR 2022.
\bibitem{attention} Vaswani et al., \emph{Attention Is All You Need}, NeurIPS 2017.
\bibitem{lamb} You et al., \emph{Large Batch Optimization for Deep Learning: Training BERT in 76 Minutes}, ICLR 2020.
\bibitem{fourier} Tancik et al., \emph{Fourier Features Let Networks Learn High Frequency Functions}, NeurIPS 2020.
\bibitem{geglu} Shazeer, \emph{GLU Variants Improve Transformer}, arXiv 2020.
\end{thebibliography}
\end{frame}

% ---------------------------------------------------------------------
\begin{frame}{Appendice: domande attese all'esame}
\small
\begin{columns}[T]
\begin{column}{0.48\textwidth}
\begin{itemize}
  \item Perché cross-attention invece di self-attention?
  \item Come cambia la complessità?
  \item A cosa serve il weight sharing?
  \item Perché Fourier PE è così importante?
\end{itemize}
\end{column}
\begin{column}{0.48\textwidth}
\begin{itemize}
  \item Perché il gap su CIFAR-10 è grande?
  \item Perché byte-level tokenization?
  \item Come migliora Perceiver IO?
  \item Può sostituire ViT/BERT?
\end{itemize}
\end{column}
\end{columns}

\begin{block}{Risposta breve da ricordare}
Perceiver usa latenti per rendere l'attention scalabile; Perceiver IO usa output queries per rendere generale anche il lato output.
\end{block}
\end{frame}

% ---------------------------------------------------------------------
\begin{frame}[plain]
\centering
\vfill
{\Huge\color{DeepBlue}\textbf{Grazie!}}

\vspace{1em}
{\Large Domande?}

\vspace{2em}
Francesco Gigli e Corso Vignoli
\vfill
\end{frame}

\end{document}
